\chapter{$U$, the Curiosity Term, and Why It Shrinks}
\label{ch:u_curiosity}

The exploration bonus for an arrow $a$ is defined as:
\begin{equation}
  U(a) = c \cdot P(a) \cdot \frac{\sqrt{N}}{1 + n(a)}
\end{equation}
where:
\begin{itemize}[leftmargin=2em]
  \item $P(a)$ = the policy prior for this arrow (fixed, from the network)
  \item $n(a)$ = visits to \textbf{THIS} arrow (grows as we use it)
  \item $N$ = total visits made from this node (grows as we use any arrow)
  \item $c$ = exploration constant (about $1.745$ in our trace)
\end{itemize}

\begin{established}[Read Each Piece Physically]
\begin{itemize}[leftmargin=2em]
  \item \textbf{$P(a)$ in the numerator} --- the network's hunch gets a say. A move it likes starts loud.
  \item \textbf{$n(a)$ in the denominator} --- the more we have already looked down this arrow, the less curiosity it deserves. \textbf{Visit an arrow once and its $U$ halves} (the denominator goes $1 \to 2$). Visit it again and it drops to a third. This is the mechanism that forces the search to move on.
  \item \textbf{$\sqrt{N}$ in the numerator} --- as the node as a whole gets more attention, everything gets a mild lift, so a promising-but-neglected arrow can come back into contention later. The square root makes this grow slowly.
\end{itemize}
\end{established}

So $U$ is large when a move is well-liked and under-examined, and small when it is either disliked or already well-examined. That is the entire idea.

\begin{center}
\begin{minipage}{\linewidth}
\centering
% fig_c_ushrink.tex -- U decay against visit count n
\begin{tikzpicture}[scale=0.82, every node/.style={font=\small}]
  % Axes
  \draw[WorkGrey, line width=0.7pt, ->] (0,0) -- (6.5, 0) node[right, font=\footnotesize] {Visits $n(a)$};
  \draw[WorkGrey, line width=0.7pt, ->] (0,0) -- (0, 4.4) node[above, font=\footnotesize] {Curiosity Bonus $U(a)$};

  % Gridlines
  \foreach \y/\val in {1.0/0.20, 2.0/0.40, 3.0/0.60, 4.0/0.80} {
    \draw[WorkGrey!30, line width=0.3pt] (0, \y) -- (6.0, \y);
    \node[left, font=\tiny, text=WorkGrey] at (0, \y) {\val};
  }
  \foreach \x in {0, 1, 2, 3, 4, 5} {
    \draw[WorkGrey, line width=0.5pt] (\x, 0.05) -- (\x, -0.05);
    \node[below, font=\footnotesize] at (\x, -0.05) {$n{=}\x$};
  }

  % High P move (Kd6: P = 0.4513)
  \draw[BuildBlue, line width=1.8pt] 
    (0, 3.938) -- (1, 1.969) -- (2, 1.313) -- (3, 0.984) -- (4, 0.788) -- (5, 0.656);

  \foreach \x/\y/\lbl in {0/3.938/0.788, 1/1.969/0.394, 2/1.313/0.263, 3/0.984/0.197, 4/0.788/0.158, 5/0.656/0.131} {
    \fill[BuildBlue] (\x, \y) circle (2.5pt);
    \node[above right=1pt, font=\tiny\bfseries, text=BuildBlue] at (\x, \y) {\lbl};
  }

  % Low P move (Kf5: P = 0.0538)
  \draw[PitfallRed, line width=1.4pt, dashed] 
    (0, 0.469) -- (1, 0.235) -- (2, 0.156) -- (3, 0.117) -- (4, 0.094) -- (5, 0.078);
  \foreach \x/\y in {0/0.469, 1/0.235, 2/0.156, 3/0.117, 4/0.094, 5/0.078} {
    \fill[PitfallRed] (\x, \y) circle (2pt);
  }
  \node[above=2pt, font=\tiny, text=PitfallRed] at (0, 0.469) {0.094};

  % Halving callout
  \draw[GoldPath, line width=1.2pt, -{Stealth[length=2.5mm]}] (1.8, 3.2) -- (0.8, 2.3);
  \node[draw=GoldPath, fill=white, rounded corners=2pt, font=\scriptsize\bfseries, text=GoldPath, align=center] at (3.2, 3.5) {%
    First Visit Halving:\\
    $\frac{1}{1+0} \to \frac{1}{1+1} \implies \mathbf{50\%\ \text{drop!}}$%
  };

  % Labels
  \node[font=\footnotesize\bfseries, text=BuildBlue, anchor=west] at (3.5, 1.8) {High-$P$ Move (Kd6, $P=45.1\%$)};
  \node[font=\footnotesize\bfseries, text=PitfallRed, anchor=west] at (3.5, 0.6) {Low-$P$ Move (Kf5, $P=5.4\%$)};
\end{tikzpicture}

\captionof{figure}{\textbf{Curiosity bonus $U$ decay as visit count $n$ grows.} The $1/(1+n)$ factor halves $U$ on the very first visit ($n = 0 \to 1$), forcing exploration of alternatives.}
\label{fig:ushrink}
\end{minipage}
\end{center}
